% =================================================================
\section{Functions, relations and graphs (Area of Study 1)}
\label{sec:functions}
% =================================================================

\begin{sdbox}
\textbf{This area of study includes:}
\begin{itemize}
  \item graphs of polynomial functions and their key features
  \item graphs of power functions $y=x^n$, $n\in\Q$; exponential functions $y=a^x$, $a\in\R^+$, in particular $y=e^x$; logarithmic functions $y=\loge(x)$ and $y=\log_{10}(x)$; and circular functions $y=\sin(x)$, $y=\cos(x)$ and $y=\tan(x)$, and their key features
  \item transformation from $y=f(x)$ to $y=Af(n(x+b))+c$, where $A,n,b,c\in\R$, $A,n\ne0$, and $f$ is one of the functions above, and the inverse transformation
  \item the relation between the graph of an original function and the graph of a transformed function (including \emph{families} of transformed functions for a single transformation parameter)
  \item graphs of sum, difference, product and composite functions of the types above (not composites that give reciprocal or quotient functions)
  \item modelling of practical situations using polynomial, power, circular, exponential and logarithmic functions, simple transformations and combinations of these, including simple piecewise (hybrid) functions
\end{itemize}
\textbf{Key skill:} \byhand\ sketch graphs of polynomials up to degree 4; $y=x^n$, $n\in\N$; $y=a^x$ (using key points $\left(-1,\frac1a\right)$, $(0,1)$, $(1,a)$); $\loge(x)$; $\log_{10}(x)$; and simple transformations of these, \emph{clearly identifying key features including any vertical or horizontal asymptotes}. Also: domain, maximal domain, co-domain, range, asymptotic behaviour, symmetry, and recognising possible models from graphs or tables.
\end{sdbox}

\subsection{Functions, domain and range}
\label{sec:domain}

\begin{defn}{Definitions}
\begin{itemize}
  \item A \textbf{relation} is a set of ordered pairs. A \textbf{function} is a relation where each $x$-value maps to \emph{exactly one} $y$-value (passes the \emph{vertical line test}).
  \item \textbf{One-to-one}: each $y$-value comes from exactly one $x$-value (passes the \emph{horizontal line test}). Otherwise \textbf{many-to-one} (e.g.\ $y=x^2$ on $\R$). A function has an inverse \emph{function} only if it is one-to-one.
  \item $f:D\to\R,\ f(x)=\text{rule}$. The \textbf{domain} $D$ is the set of allowed inputs; the \textbf{co-domain} ($\R$ in Methods) contains all outputs; the \textbf{range} is the set of outputs actually produced.
  \item The \textbf{maximal} (implied or natural) domain is the largest subset of $\R$ for which the rule is defined. When a function is given by a rule only, assume the maximal domain.
  \item \textbf{Even}: $f(-x)=f(x)$, so the graph is symmetric about the $y$-axis ($x^2$, $x^4$, $\cos x$, $\frac{1}{x^2}$). \textbf{Odd}: $f(-x)=-f(x)$, so the graph has rotational symmetry about the origin ($x^3$, $\frac1x$, $\sin x$, $\tan x$, $x^{1/3}$).
\end{itemize}
\end{defn}

\begin{cols}
\begin{keyres}{Finding the maximal domain}
\begin{itemize}
  \item Denominator $\ne0$.
  \item Even root $\sqrt{\ \ }$: inside $\ge0$.
  \item $\loge(\ \ )$: inside $>0$.
  \item $\tan(\theta)$: $\theta\ne\frac{\pi}{2}+k\pi$, $k\in\Z$.
  \item Sum, difference, product: \emph{intersection} of the domains, $d_{f\pm g}=d_{fg}=d_f\cap d_g$.
\end{itemize}
\end{keyres}
\nextcol
\begin{keyres}{Finding the range}
\begin{itemize}
  \item Sketch it. Read off from the lowest to the highest $y$-value, using endpoints, turning points and asymptotes.
  \item Check endpoints: closed ($\bullet$, included) or open ($\circ$, excluded).
  \item Asymptotes are \emph{never} included in the range ($y=e^x$ has range $(0,\infty)$).
\end{itemize}
\end{keyres}
\end{cols}

\begin{example}{maximal domain and range \byhand}
\begin{tabularx}{\linewidth}{@{}p{3.8cm}L@{}}
$f(x)=\sqrt{4-2x}+1$ & Need $4-2x\ge0\Rightarrow x\le2$. Domain $(-\infty,2]$; range $[1,\infty)$ (endpoint $(2,1)$). \\
$g(x)=\loge(3-x)-2$ & Need $3-x>0\Rightarrow x<3$. Domain $(-\infty,3)$; range $\R$; vertical asymptote $x=3$; $x$-intercept $3-e^2$; $y$-intercept $\loge3-2$. \\
$h(x)=\dfrac{2}{(x-1)^2}-3$ & Domain $\R\setminus\{1\}$; asymptotes $x=1$, $y=-3$; range $(-3,\infty)$. \\
\end{tabularx}
\end{example}

\begin{tip}[Key features checklist for any sketch]
Axis intercepts (as coordinates) \textbullet\ stationary points and their nature \textbullet\ points of inflection (if asked) \textbullet\ asymptotes as dashed lines \emph{with equations} \textbullet\ endpoints shown closed $\bullet$ or open $\circ$ with coordinates \textbullet\ correct behaviour near asymptotes and as $x\to\pm\infty$ \textbullet\ correct curvature (don't draw a parabola-shaped log graph) \textbullet\ symmetry used where possible.
\end{tip}

% -----------------------------------------------------------------
\subsection{Polynomial functions}
\label{sec:poly}

\begin{keyres}{Shape rules for polynomials of degree $n$}
\begin{itemize}
  \item At most $n$ $x$-intercepts and at most $n-1$ turning points. The $y$-intercept is the constant term.
  \item \textbf{Behaviour at a root} $x=a$ from a factor $(x-a)^k$: $k=1$ the graph \emph{crosses}; $k=2$ it \emph{touches} (turning point on the axis); $k=3$ it crosses with a \emph{stationary point of inflection} (flattens).
  \item \textbf{End behaviour} (leading term $ax^n$ dominates): $n$ even, $a>0$: up at both ends; $n$ even, $a<0$: down at both ends; $n$ odd, $a>0$: bottom left to top right; $n$ odd, $a<0$: top left to bottom right.
  \item Turning-point (stationary-point) forms: $y=a(x-h)^2+k$ and $y=a(x-h)^4+k$ have a turning point at $(h,k)$; $y=a(x-h)^3+k$ has a stationary point of inflection at $(h,k)$.
  \item Comparing powers: for $|x|<1$, $x^4<x^2$ (flatter near $0$); for $|x|>1$, $x^4>x^2$. All of $x^2,x^4$ pass through $(-1,1),(0,0),(1,1)$; $x^3,x^5$ through $(-1,-1),(0,0),(1,1)$.
\end{itemize}
\end{keyres}

\begin{cols}
\begin{example}{a cubic \byhand}
Sketch $y=(x+1)^2(x-2)$.\par
\textbf{Intercepts:} touches at $x=-1$, crosses at $x=2$; $y$-intercept $(1)^2(-2)=-2$.\par
\textbf{Stationary points:} $\dydx=2(x+1)(x-2)+(x+1)^2=3(x+1)(x-1)$, so local max $(-1,0)$ and local min $(1,-4)$.\par
\textbf{Ends:} degree 3, positive leading coefficient: bottom left to top right.
\end{example}
\nextcol
\centering
\begin{tikzpicture}
\begin{axis}[mygraph, width=6.4cm, height=5.0cm, xmin=-2.3, xmax=3.0, ymin=-5.3, ymax=3.4, xtick=\empty, ytick=\empty]
  \addplot[plotA, domain=-1.95:2.45] {(x+1)^2*(x-2)};
  \node[pt] at (axis cs:-1,0) {}; \node[lab, above left] at (axis cs:-1,0) {$(-1,0)$};
  \node[pt] at (axis cs:1,-4) {}; \node[lab, below] at (axis cs:1,-4) {$(1,-4)$};
  \node[pt] at (axis cs:2,0) {};  \node[lab, below right] at (axis cs:2,0) {$(2,0)$};
  \node[pt] at (axis cs:0,-2) {}; \node[lab, left] at (axis cs:0,-2) {$(0,-2)$};
\end{axis}
\end{tikzpicture}
\end{cols}

\begin{cols}
\begin{example}{a quartic \byhand}
Sketch $y=x(x-2)^3$.\par
\textbf{Intercepts:} crosses at $(0,0)$; at $x=2$ (triple root) a stationary point of inflection.\par
\textbf{Stationary points:} $\dydx=(x-2)^3+3x(x-2)^2=(x-2)^2(4x-2)$: $x=\frac12$ gives the local min $\left(\frac12,-\frac{27}{16}\right)$; $x=2$ gives the stationary inflection $(2,0)$.\par
\textbf{Ends:} degree 4, $a>0$: up at both ends.
\end{example}
\nextcol
\centering
\begin{tikzpicture}
\begin{axis}[mygraph, width=6.4cm, height=5.0cm, xmin=-0.9, xmax=3.0, ymin=-2.4, ymax=4.2, xtick=\empty, ytick=\empty]
  \addplot[plotB, domain=-0.42:2.9] {x*(x-2)^3};
  \node[pt] at (axis cs:0,0) {}; \node[lab, above left] at (axis cs:0,0) {$(0,0)$};
  \node[pt] at (axis cs:0.5,-1.6875) {}; \node[lab, below] at (axis cs:0.5,-1.6875) {$\left(\frac12,-\frac{27}{16}\right)$};
  \node[pt] at (axis cs:2,0) {}; \node[lab, below right] at (axis cs:2,0) {$(2,0)$ stat.\ inflection};
\end{axis}
\end{tikzpicture}
\end{cols}

% -----------------------------------------------------------------
\subsection{Power functions $y=x^n$, $n\in\Q$}
\label{sec:power}

\begin{center}
\begin{tikzpicture}
\begin{axis}[smallgraph, title={\small $n\in\N$ even: $x^2$, $x^4$}, xmin=-1.6, xmax=1.6, ymin=-0.5, ymax=2.2, xtick=\empty, ytick=\empty]
  \addplot[plotA, domain=-1.45:1.45] {x^2};
  \addplot[plotB, domain=-1.2:1.2] {x^4};
  \node[pt] at (axis cs:1,1) {}; \node[lab, right] at (axis cs:1,1) {$(1,1)$};
  \node[pt] at (axis cs:-1,1) {}; \node[lab, left] at (axis cs:-1,1) {$(-1,1)$};
\end{axis}
\end{tikzpicture}\hfill
\begin{tikzpicture}
\begin{axis}[smallgraph, title={\small $n\in\N$ odd: $x^3$, $x^5$}, xmin=-1.6, xmax=1.6, ymin=-2, ymax=2, xtick=\empty, ytick=\empty]
  \addplot[plotA, domain=-1.25:1.25] {x^3};
  \addplot[plotB, domain=-1.14:1.14] {x^5};
  \node[pt] at (axis cs:1,1) {}; \node[lab, right] at (axis cs:1,1) {$(1,1)$};
  \node[pt] at (axis cs:-1,-1) {}; \node[lab, left] at (axis cs:-1,-1) {$(-1,-1)$};
\end{axis}
\end{tikzpicture}\hfill
\begin{tikzpicture}
\begin{axis}[smallgraph, title={\small $n=-1$: $y=\frac1x$}, xmin=-3, xmax=3, ymin=-3, ymax=3, xtick=\empty, ytick=\empty]
  \addplot[plotA, domain=-3:-0.34] {1/x};
  \addplot[plotA, domain=0.34:3] {1/x};
  \node[pt] at (axis cs:1,1) {}; \node[lab, above right] at (axis cs:1,1) {$(1,1)$};
  \node[pt] at (axis cs:-1,-1) {}; \node[lab, below left] at (axis cs:-1,-1) {$(-1,-1)$};
\end{axis}
\end{tikzpicture}

\begin{tikzpicture}
\begin{axis}[smallgraph, title={\small $n=-2$: $y=\frac{1}{x^2}$}, xmin=-3, xmax=3, ymin=-0.6, ymax=3.2, xtick=\empty, ytick=\empty]
  \addplot[plotA, domain=-3:-0.57] {1/x^2};
  \addplot[plotA, domain=0.57:3] {1/x^2};
  \node[pt] at (axis cs:1,1) {}; \node[lab, above right] at (axis cs:1,1) {$(1,1)$};
  \node[pt] at (axis cs:-1,1) {}; \node[lab, above left] at (axis cs:-1,1) {$(-1,1)$};
\end{axis}
\end{tikzpicture}\hfill
\begin{tikzpicture}
\begin{axis}[smallgraph, title={\small $x^{1/2}=\sqrt{x}$,\ \ $x^{1/3}=\sqrt[3]{x}$}, xmin=-2.6, xmax=2.6, ymin=-1.6, ymax=1.8, xtick=\empty, ytick=\empty]
  \addplot[plotA, domain=0:2.5, samples=300] {sqrt(x)};
  \addplot[plotB, dashed, domain=-2.5:2.5] {sign(x)*abs(x)^(1/3)};
  \node[pt] at (axis cs:1,1) {}; \node[lab, above] at (axis cs:1,1) {$(1,1)$};
  \node[pt] at (axis cs:-1,-1) {}; \node[lab, below] at (axis cs:-1,-1) {$(-1,-1)$};
\end{axis}
\end{tikzpicture}\hfill
\begin{tikzpicture}
\begin{axis}[smallgraph, title={\small $x^{2/3}$ (cusp at $O$)}, xmin=-2.6, xmax=2.6, ymin=-0.4, ymax=2.0, xtick=\empty, ytick=\empty]
  \addplot[plotA, domain=-2.5:2.5] {abs(x)^(2/3)};
  \node[pt] at (axis cs:1,1) {}; \node[lab, above] at (axis cs:1,1) {$(1,1)$};
  \node[pt] at (axis cs:-1,1) {}; \node[lab, above] at (axis cs:-1,1) {$(-1,1)$};
\end{axis}
\end{tikzpicture}
\end{center}

\begin{center}
\small
\begin{tabularx}{\linewidth}{@{}lllL@{}}
\toprule
$y=x^n$ & Maximal domain & Range & Notes \\
\midrule
$n\in\N$ even ($x^2,x^4$) & $\R$ & $[0,\infty)$ & even; turning point at $O$ \\
$n\in\N$ odd ($x^3,x^5$) & $\R$ & $\R$ & odd; stationary inflection at $O$ ($n\ge3$) \\
$n=-1$, $\frac1x$ (hyperbola) & $\R\setminus\{0\}$ & $\R\setminus\{0\}$ & odd; asymptotes $x=0$, $y=0$ \\
$n=-2$, $\frac{1}{x^2}$ (truncus) & $\R\setminus\{0\}$ & $(0,\infty)$ & even; asymptotes $x=0$, $y=0$ \\
$n=\frac12$, $\sqrt{x}$ & $[0,\infty)$ & $[0,\infty)$ & starts at $(0,0)$; vertical tangent there \\
$n=\frac13$, $\sqrt[3]{x}$ & $\R$ & $\R$ & odd; vertical tangent at $O$ (not differentiable at $0$) \\
$n=\frac23$, $\left(\sqrt[3]{x}\right)^2$ & $\R$ & $[0,\infty)$ & even; cusp at $O$ \\
$n=-\frac12$, $\frac{1}{\sqrt{x}}$ & $(0,\infty)$ & $(0,\infty)$ & asymptotes $x=0$, $y=0$ \\
\bottomrule
\end{tabularx}
\end{center}

\begin{keyres}{General rational powers and transformed power functions}
For $n=\frac{p}{q}$ in simplest form ($q>0$): $x^{p/q}=\left(\sqrt[q]{x}\right)^p$. If $q$ is \textbf{odd}, negative $x$ is allowed; if $q$ is \textbf{even}, need $x\ge0$. If $p<0$, also exclude $x=0$. Every $y=x^n$ passes through $(1,1)$.
\begin{itemize}
  \item $y=\dfrac{a}{x-h}+k$ and $y=\dfrac{a}{(x-h)^2}+k$: asymptotes $x=h$ and $y=k$.
  \item $y=a\sqrt{x-h}+k$: starts at $(h,k)$, domain $[h,\infty)$; $y=a\sqrt{h-x}+k$ has domain $(-\infty,h]$.
\end{itemize}
\end{keyres}

% -----------------------------------------------------------------
\subsection{Exponential and logarithmic functions}
\label{sec:explog}

\begin{cols}
\begin{tikzpicture}
\begin{axis}[mygraph, width=6.4cm, height=5.2cm, xmin=-2.6, xmax=2.6, ymin=-0.6, ymax=5.2, xtick=\empty, ytick=\empty]
  \addplot[plotA, domain=-2.5:2.3] {2^x};
  \addplot[plotC, domain=-2.5:1.62] {exp(x)};
  \addplot[plotB, domain=-2.3:2.5] {0.5^x};
  \node[pt] at (axis cs:0,1) {}; \node[lab, above left] at (axis cs:0,1) {$(0,1)$};
  \node[pt] at (axis cs:1,2) {}; \node[lab, right] at (axis cs:1,2) {$(1,2)$};
  \node[pt] at (axis cs:-1,0.5) {}; \node[lab, below] at (axis cs:-1,0.5) {$\left(-1,\frac12\right)$};
  \node[lab, plotA] at (axis cs:2.05,4.7) {$2^x$};
  \node[lab, plotC] at (axis cs:1.05,4.4) {$e^x$};
  \node[lab, plotB] at (axis cs:-2.0,4.6) {$\left(\frac12\right)^x$};
  \node[lab] at (axis cs:2.1,0.25) {$y=0$};
\end{axis}
\end{tikzpicture}
\nextcol
\begin{keyres}{$y=a^x$, $a>0$, $a\ne1$}
\begin{itemize}
  \item Domain $\R$; range $(0,\infty)$; horizontal asymptote $y=0$.
  \item Key points $\left(-1,\frac1a\right)$, $(0,1)$, $(1,a)$.
  \item $a>1$: strictly increasing (growth). $0<a<1$: strictly decreasing (decay). $\left(\frac1a\right)^x=a^{-x}$ is the reflection in the $y$-axis.
  \item $e\approx2.718$. $y=e^x$ has gradient $1$ at $(0,1)$ and $\ddx e^x=e^x$.
  \item $y=Ae^{n(x-h)}+k$: asymptote $y=k$; range $(k,\infty)$ if $A>0$, $(-\infty,k)$ if $A<0$.
\end{itemize}
\end{keyres}
\end{cols}

\begin{cols}
\begin{keyres}{$y=\log_a(x)$: the inverse of $y=a^x$}
\begin{itemize}
  \item $y=\log_a(x)\iff a^y=x$ (for $x>0$).
  \item Domain $(0,\infty)$; range $\R$; vertical asymptote $x=0$.
  \item Key points $\left(\frac1a,-1\right)$, $(1,0)$, $(a,1)$; reflection of $y=a^x$ in $y=x$.
  \item $\loge(x)=\ln(x)$ passes through $(e,1)$; $\log_{10}(x)$ through $(10,1)$, flatter than $\loge$ for $x>1$.
  \item $y=A\loge(n(x-h))+k$: asymptote $x=h$; domain $x>h$ if $n>0$, $x<h$ if $n<0$.
\end{itemize}
\end{keyres}
\nextcol
\centering
\begin{tikzpicture}
\begin{axis}[mygraph, width=6.4cm, height=5.6cm, xmin=-2.6, xmax=3.6, ymin=-2.6, ymax=3.6, xtick=\empty, ytick=\empty]
  \addplot[black!40, dashed, thin, domain=-2.5:3.5] {x};
  \addplot[plotC, domain=-2.5:1.2] {exp(x)};
  \addplot[plotA, domain=0.07:3.5, samples=300] {ln(x)};
  \addplot[plotB, domain=0.06:3.5, samples=300] {log10(x)};
  \node[pt] at (axis cs:1,0) {}; \node[lab, below right] at (axis cs:1,0) {$(1,0)$};
  \node[pt] at (axis cs:2.71828,1) {}; \node[lab, above left] at (axis cs:2.71828,1) {$(e,1)$};
  \node[pt] at (axis cs:0,1) {}; \node[lab, left] at (axis cs:0,1) {$(0,1)$};
  \node[lab, plotC] at (axis cs:1.35,3.3) {$e^x$};
  \node[lab, plotA] at (axis cs:3.2,1.45) {$\loge x$};
  \node[lab, plotB] at (axis cs:3.2,0.25) {$\log_{10}x$};
  \node[lab, black!60] at (axis cs:2.9,3.35) {$y=x$};
\end{axis}
\end{tikzpicture}
\end{cols}

\begin{keyres}{Index laws and logarithm laws (memorise)}
\begin{cols}
\textbf{Index laws} ($a,b>0$)
\begin{itemize}
  \item $a^ma^n=a^{m+n}$, \ $\dfrac{a^m}{a^n}=a^{m-n}$, \ $(a^m)^n=a^{mn}$
  \item $(ab)^n=a^nb^n$, \ $a^0=1$, \ $a^{-n}=\dfrac{1}{a^n}$
  \item $a^{p/q}=\sqrt[q]{a^p}=\left(\sqrt[q]{a}\right)^p$
  \item $a^{x}=e^{x\loge a}$ (useful for calculus)
\end{itemize}
\nextcol
\textbf{Logarithm laws} ($a>0$, $a\ne1$, $m,n>0$)
\begin{itemize}
  \item $\log_a(mn)=\log_a m+\log_a n$; \ $\log_a\!\left(\frac{m}{n}\right)=\log_a m-\log_a n$
  \item $\log_a(m^p)=p\log_a m$; \ $\log_a\!\left(\frac1m\right)=-\log_a m$
  \item $\log_a1=0$, \ $\log_a a=1$, \ $a^{\log_a x}=x$, \ $\log_a(a^x)=x$
  \item Change of base: $\log_a x=\dfrac{\loge x}{\loge a}$
\end{itemize}
\end{cols}
\end{keyres}

\begin{warn}
$\log_a(m+n)\ne\log_a m+\log_a n$ \textbullet\ $\dfrac{\log m}{\log n}\ne\log\!\left(\dfrac mn\right)$ \textbullet\ $(\loge x)^2\ne2\loge x$ \textbullet\ $\loge(x^2)=2\loge|x|$, so $\loge(x^2)$ and $2\loge x$ have \emph{different domains} \textbullet\ $e^{a+b}=e^ae^b$, not $e^a+e^b$.
\end{warn}

\begin{example}{sketching transformed exponential and log graphs \byhand}
\textbf{(a)} $y=3-e^{2x}$: asymptote $y=3$ (as $x\to-\infty$); $y$-intercept $3-1=2$; $x$-intercept $e^{2x}=3\Rightarrow x=\frac12\loge3$; range $(-\infty,3)$; strictly decreasing.\par
\textbf{(b)} $y=\loge(3-x)-2$ (from the earlier domain example): reflect $y=\loge x$ in the $y$-axis, translate 3 right, 2 down. Asymptote $x=3$; $x$-intercept $3-e^2\approx-4.39$; $y$-intercept $\loge3-2\approx-0.90$; domain $(-\infty,3)$; strictly decreasing.
\end{example}

% -----------------------------------------------------------------
\subsection{Circular functions}
\label{sec:circ}

\begin{cols}
\begin{tikzpicture}
\begin{axis}[mygraph, width=6.6cm, height=4.3cm, xmin=-0.4, xmax=7.1, ymin=-1.5, ymax=1.6,
  xtick={1.5708,3.14159,4.71239,6.28319}, xticklabels={$\frac{\pi}{2}$,$\pi$,$\frac{3\pi}{2}$,$2\pi$},
  ytick={-1,1}, every tick label/.append style={font=\scriptsize}]
  \addplot[plotA, domain=0:6.28319] {sin(deg(x))};
  \addplot[plotB, dashed, domain=0:6.28319] {cos(deg(x))};
  \node[lab, plotA] at (axis cs:1.0,1.25) {$\sin x$};
  \node[lab, plotB] at (axis cs:5.6,1.25) {$\cos x$};
\end{axis}
\end{tikzpicture}

\begin{tikzpicture}
\begin{axis}[mygraph, width=6.6cm, height=4.6cm, xmin=-3.4, xmax=3.6, ymin=-4, ymax=4,
  xtick={-3.14159,-1.5708,1.5708,3.14159}, xticklabels={$-\pi$,$-\frac{\pi}{2}$,$\frac{\pi}{2}$,$\pi$}, ytick=\empty,
  restrict y to domain=-6:6, unbounded coords=jump]
  \addplot[plotA, domain=-3.14159:3.14159, samples=401] {tan(deg(x))};
  \asym (axis cs:-1.5708,-4) -- (axis cs:-1.5708,4);
  \asym (axis cs:1.5708,-4) -- (axis cs:1.5708,4);
  \node[lab, plotA] at (axis cs:2.6,2.6) {$\tan x$};
\end{axis}
\end{tikzpicture}
\nextcol
\begin{keyres}{Key features}
\begin{itemize}
  \item $\sin$, $\cos$: domain $\R$, range $[-1,1]$, period $2\pi$. $\cos x=\sin\!\left(x+\frac{\pi}{2}\right)$.
  \item $\tan$: period $\pi$; asymptotes $x=\frac{\pi}{2}+k\pi$; zeros at $x=k\pi$; range $\R$.
  \item $y=a\sin(n(x-b))+c$ or $a\cos(n(x-b))+c$:\\
   amplitude $|a|$; \ period $\dfrac{2\pi}{|n|}$; \ mean line $y=c$; \ range $[c-|a|,\,c+|a|]$; \ horizontal shift $b$.
  \item $y=a\tan(n(x-b))+c$: period $\dfrac{\pi}{|n|}$; asymptotes where $n(x-b)=\frac{\pi}{2}+k\pi$, i.e.\ $x=b+\dfrac{\pi}{2n}+\dfrac{k\pi}{n}$, $k\in\Z$.
\end{itemize}
\end{keyres}
\begin{method}{Sketching $y=a\sin(n(x-b))+c$ \byhand}
\begin{enumerate}
  \item Factorise the inside: $\sin(2x-\frac{\pi}{2})=\sin\!\left(2\left(x-\frac{\pi}{4}\right)\right)$.
  \item Period, amplitude, mean line, range.
  \item Mark a cycle starting at $x=b$ and split it into quarters. The key points are mean, max, mean, min, mean (sin) or max, mean, min, mean, max (cos), flipped if $a<0$.
  \item Find the intercepts by solving; plot the endpoints of the domain.
\end{enumerate}
\end{method}
\end{cols}

\begin{cols}[0.52]
\begin{example}{a transformed circular graph \byhand}
Sketch $y=2\sin\!\left(2x-\frac{\pi}{2}\right)+1$ for $x\in[0,2\pi]$.\par
$y=2\sin\!\left(2\left(x-\frac{\pi}{4}\right)\right)+1$: amplitude $2$, period $\pi$, mean line $y=1$, range $[-1,3]$. It simplifies to $y=1-2\cos(2x)$.\par
$y$-intercept: $2\sin\!\left(-\frac{\pi}{2}\right)+1=-1$.\par
$x$-intercepts: $\cos(2x)=\frac12$, $2x\in[0,4\pi]$, so $2x=\frac{\pi}{3},\frac{5\pi}{3},\frac{7\pi}{3},\frac{11\pi}{3}$, giving $x=\frac{\pi}{6},\frac{5\pi}{6},\frac{7\pi}{6},\frac{11\pi}{6}$.\par
Maxima $\left(\frac{\pi}{2},3\right)$, $\left(\frac{3\pi}{2},3\right)$; minima $(0,-1)$, $(\pi,-1)$, $(2\pi,-1)$.
\end{example}
\nextcol[0.455]
\centering
\begin{tikzpicture}
\begin{axis}[mygraph, width=7.0cm, height=5.0cm, xmin=-0.3, xmax=6.9, ymin=-1.6, ymax=3.6,
  xtick={0.5236,1.5708,3.14159,4.71239,6.28319}, xticklabels={$\frac{\pi}{6}$,$\frac{\pi}{2}$,$\pi$,$\frac{3\pi}{2}$,$2\pi$},
  ytick={-1,1,3}, every tick label/.append style={font=\scriptsize}]
  \addplot[black!40, dashed, thin, domain=0:6.28319] {1};
  \addplot[plotA, domain=0:6.28319] {1-2*cos(deg(2*x))};
  \node[pt] at (axis cs:0,-1) {}; \node[pt] at (axis cs:6.28319,-1) {};
  \node[pt] at (axis cs:1.5708,3) {}; \node[lab, above] at (axis cs:1.5708,3) {$\left(\frac{\pi}{2},3\right)$};
  \node[pt] at (axis cs:4.71239,3) {}; \node[lab, above] at (axis cs:4.71239,3) {$\left(\frac{3\pi}{2},3\right)$};
  \node[pt] at (axis cs:3.14159,-1) {}; \node[lab, below] at (axis cs:3.14159,-1) {$(\pi,-1)$};
  \node[pt] at (axis cs:0.5236,0) {}; \node[pt] at (axis cs:2.618,0) {};
  \node[pt] at (axis cs:3.6652,0) {}; \node[pt] at (axis cs:5.7596,0) {};
  \node[lab, black!60] at (axis cs:6.55,1.25) {$y=1$};
\end{axis}
\end{tikzpicture}
\end{cols}

% -----------------------------------------------------------------
\subsection{Transformations of graphs}
\label{sec:trans}

\begin{keyres}{The study design form $y=Af\big(n(x+b)\big)+c$, \ $A,n\ne0$}
\begin{center}
\begin{tabularx}{\linewidth}{@{}lL@{}}
\toprule
Parameter & Effect on the graph of $y=f(x)$ \\
\midrule
$|A|$ & Dilation by factor $|A|$ \textbf{from the $x$-axis} (parallel to the $y$-axis). If $A<0$: also a \textbf{reflection in the $x$-axis}. \\
$|n|$ & Dilation by factor $\dfrac{1}{|n|}$ \textbf{from the $y$-axis} (parallel to the $x$-axis). If $n<0$: also a \textbf{reflection in the $y$-axis}. \\
$b$ & Translation of $b$ units in the \textbf{negative} $x$-direction ($x+3$: 3 units \emph{left}; $x-3$: 3 units \emph{right}). \\
$c$ & Translation of $c$ units in the positive $y$-direction (up if $c>0$). \\
\bottomrule
\end{tabularx}
\end{center}
\textbf{Mapping of points:} \quad $(x,\,y)\ \longrightarrow\ \left(\dfrac{x}{n}-b,\ \ Ay+c\right)$. \quad For this form apply \textbf{dilations and reflections first, then translations}.\par
\textbf{Effect on features:} the image of a point is found with the mapping. A vertical asymptote $x=h$ becomes $x=\frac{h}{n}-b$; a horizontal asymptote $y=k$ becomes $y=Ak+c$; maxima become minima under a reflection in the $x$-axis; domain and range transform with the $x$- and $y$-rules respectively.
\end{keyres}

\begin{method}{Finding the transformations between two rules}
\textbf{Method 1 (read off):} rewrite the image rule as $y=Af(n(x+b))+c$, \emph{taking out the factor $n$ inside the brackets first}, then read off $A$, $n$, $b$, $c$.\par
\textbf{Method 2 (dash method):} let the image of $(x,y)$ be $(x',y')$. Write $x'$ and $y'$ in terms of $x$ and $y$, make $x$ and $y$ the subjects, and substitute into $y=f(x)$. Then drop the dashes. This works for \emph{any} described sequence of transformations, in any order.
\end{method}

\begin{cols}[0.6]
\begin{example}{describe the transformations \byhand}
Find a sequence of transformations that maps $y=\sqrt{x}$ onto $y=3-2\sqrt{4-x}$.\par
Write it as $y=-2\sqrt{-(x-4)}+3$: $A=-2$, $n=-1$, $b=-4$, $c=3$.
\begin{enumerate}
  \item dilation by factor 2 from the $x$-axis;
  \item reflection in the $x$-axis;
  \item reflection in the $y$-axis;
  \item translation 4 units right and 3 units up.
\end{enumerate}
Mapping: $(x,y)\to(4-x,\ 3-2y)$.\par
\textbf{Check (dash method):} $x'=4-x\Rightarrow x=4-x'$, \ $y'=3-2y\Rightarrow y=\frac{3-y'}{2}$. Substitute into $y=\sqrt x$: $\frac{3-y'}{2}=\sqrt{4-x'}$, so $y'=3-2\sqrt{4-x'}$. \tick
\end{example}
\nextcol[0.37]
\centering
\begin{tikzpicture}
\begin{axis}[mygraph, width=5.6cm, height=5.6cm, xmin=-1.2, xmax=5.4, ymin=-2.2, ymax=3.8, xtick=\empty, ytick=\empty]
  \addplot[plotA, domain=0:5, samples=300] {sqrt(x)};
  \addplot[plotB, domain=-1:4, samples=300] {3-2*sqrt(4-x)};
  \node[pt] at (axis cs:0,0) {};
  \node[pt] at (axis cs:4,3) {}; \node[lab, above] at (axis cs:4,3) {$(4,3)$};
  \node[pt] at (axis cs:1.75,0) {}; \node[lab, below right] at (axis cs:1.75,0) {$\left(\frac74,0\right)$};
  \node[pt] at (axis cs:0,-1) {}; \node[lab, left] at (axis cs:0,-1) {$(0,-1)$};
  \node[lab, plotA] at (axis cs:4.8,1.75) {$\sqrt{x}$};
  \node[lab, plotB] at (axis cs:2.2,1.6) {image};
\end{axis}
\end{tikzpicture}

{\scriptsize The image starts at $(4,3)$ (the image of $(0,0)$) and has domain $(-\infty,4]$ and range $(-\infty,3]$.}
\end{cols}

\begin{example}{image under a mapping; non-uniqueness}
Find the image of $y=\loge(x)$ under $T(x,y)=\left(\frac{x}{3}+2,\ 2y+1\right)$.\par
$x'=\frac x3+2\Rightarrow x=3(x'-2)$; \ $y'=2y+1\Rightarrow y=\frac{y'-1}{2}$. So $\frac{y'-1}{2}=\loge(3(x'-2))$, giving $y=2\loge(3x-6)+1$.\par
\textbf{Non-uniqueness:} by the log laws, $2\loge(3(x-2))+1=2\loge(x-2)+(1+2\loge3)$. So the same graph is obtained \emph{without} any dilation from the $y$-axis: dilate by factor 2 from the $x$-axis, then translate 2 units right and $1+2\loge3$ units up. Exam questions may have several correct answers.\par
\textbf{Inverse transformation:} undo in reverse order. $T^{-1}(x,y)=\left(3(x-2),\ \frac{y-1}{2}\right)$ maps $y=2\loge(3x-6)+1$ back onto $y=\loge x$.
\end{example}

\begin{warn}
\begin{itemize}
  \item \textbf{Order matters.} For $y=x^2$: ``dilate by factor 2 from the $y$-axis, then translate 3 right'' gives $y=\left(\frac{x-3}{2}\right)^2$; ``translate 3 right, then dilate by factor 2 from the $y$-axis'' gives $y=\left(\frac{x}{2}-3\right)^2$. These are different graphs.
  \item $y=f(2x)$ is a dilation by factor $\frac12$ (not 2) from the $y$-axis. In $y=\sin(2x+\pi)=\sin\!\big(2\left(x+\frac{\pi}{2}\right)\big)$ the translation is $\frac{\pi}{2}$ left, not $\pi$.
  \item ``From the $x$-axis'' means vertical stretching. ``Parallel to the $x$-axis'' means horizontal. Read carefully.
\end{itemize}
\end{warn}

\begin{cas}[Transformations seen in other forms]
Exams sometimes write a transformation as a matrix equation:
$T\!\left(\begin{bmatrix}x\\y\end{bmatrix}\right)=\begin{bmatrix}\frac1n&0\\0&A\end{bmatrix}\begin{bmatrix}x\\y\end{bmatrix}+\begin{bmatrix}-b\\c\end{bmatrix}$, which is the same mapping $(x,y)\to\left(\frac xn-b,\ Ay+c\right)$. Read it as ``$x'=\frac{x}{n}-b$, $y'=Ay+c$'' and use the dash method.
\end{cas}

\begin{cols}[0.58]
\begin{keyres}{Families of functions (one parameter)}\label{sec:families}
A \textbf{family} is a set of graphs given by one rule with a parameter, e.g.\ $y=e^{kx}$, $y=kx(x-2)$, $y=x^3-kx$, $y=(x-k)^2+k$.
\begin{itemize}
  \item \textbf{Common points:} find points that do not depend on the parameter. $y=kx(x-2)$ always passes through $(0,0)$ and $(2,0)$; $y=e^{kx}$ through $(0,1)$; $y=k(x-1)+2$ through $(1,2)$.
  \item \textbf{Locus of a feature:} $y=(x-k)^2+k$ has turning point $(k,k)$, which always lies on $y=x$ (see the figure).
  \item \textbf{Conditions on the parameter:} e.g.\ $y=x^3-kx$ has two stationary points only if $k>0$ (they are at $x=\pm\sqrt{k/3}$).
  \item Use CAS to graph several members (Exam 2), then prove the result algebraically.
\end{itemize}
\end{keyres}
\nextcol[0.39]
\centering
\begin{tikzpicture}
\begin{axis}[mygraph, width=6.0cm, height=5.5cm, xmin=-2.6, xmax=3.6, ymin=-1.8, ymax=4.2, xtick=\empty, ytick=\empty]
  \addplot[black!40, dashed, thin, domain=-2.5:3.5] {x};
  \addplot[plotA, domain=-2.9:0.9] {(x+1)^2-1};
  \addplot[plotB, domain=-1.9:1.9] {x^2};
  \addplot[plotC, domain=-0.9:2.9] {(x-1)^2+1};
  \addplot[plotD, domain=0.1:3.9] {(x-2)^2+2};
  \node[pt] at (axis cs:-1,-1) {}; \node[pt] at (axis cs:0,0) {};
  \node[pt] at (axis cs:1,1) {};   \node[pt] at (axis cs:2,2) {};
  \node[lab, black!60] at (axis cs:3.15,3.0) {$y=x$};
\end{axis}
\end{tikzpicture}

{\scriptsize $y=(x-k)^2+k$ for $k=-1,0,1,2$.}
\end{cols}

% -----------------------------------------------------------------
\subsection{Sums, differences, products and composites of functions}
\label{sec:combo}

\begin{keyres}{Combining functions}
\begin{itemize}
  \item $(f\pm g)(x)=f(x)\pm g(x)$ and $(fg)(x)=f(x)\,g(x)$, each with domain $d_f\cap d_g$.
  \item \textbf{Addition of ordinates} (sketching $f+g$): sketch $f$ and $g$ on the same axes, then add $y$-values at key $x$-values: where either is $0$, where they are equal or opposite, and near asymptotes. Near a vertical asymptote of one function, the sum follows that asymptote; for large $|x|$ the dominant term decides the behaviour.
  \item \textbf{Products:} $y=f(x)g(x)$ is $0$ wherever either factor is $0$. The sign of the product follows from the signs of the factors. $e^{-x}\sin x$ oscillates between the curves $\pm e^{-x}$. (See \S\ref{sec:sketch} for a full sketch of $y=xe^{-x}$.)
  \item \textbf{Composites} $y=f(g(x))$: work out the range of the inner function, then apply the outer one. E.g.\ $y=e^{\sin x}$ has range $\left[e^{-1},e\right]$; $y=\loge(x^2+1)$ is even with range $[0,\infty)$. (Reciprocal or quotient composites such as $\frac{1}{\sin x}$ are \emph{not} required.)
\end{itemize}
\end{keyres}

\begin{cols}
\begin{example}{addition of ordinates \byhand}
Sketch $y=x^2+\dfrac{2}{x}$.\par
Domain $\R\setminus\{0\}$; vertical asymptote $x=0$. For large $|x|$ the graph follows $y=x^2$.\par
$x$-intercept: $x^3=-2\Rightarrow x=-\sqrt[3]{2}\approx-1.26$.\par
Stationary point: $2x-\frac{2}{x^2}=0\Rightarrow x^3=1$, so $(1,3)$, a local minimum.
\end{example}
\nextcol
\centering
\begin{tikzpicture}
\begin{axis}[mygraph, width=6.3cm, height=5.0cm, xmin=-2.6, xmax=2.6, ymin=-4.5, ymax=6.5, xtick=\empty, ytick=\empty]
  \addplot[black!45, dashed, thin, domain=-2.5:2.5] {x^2};
  \addplot[black!45, dotted, thick, domain=-2.5:-0.31] {2/x};
  \addplot[black!45, dotted, thick, domain=0.31:2.5] {2/x};
  \addplot[plotA, domain=-2.5:-0.36, samples=300] {x^2+2/x};
  \addplot[plotA, domain=0.33:2.5, samples=300] {x^2+2/x};
  \node[pt] at (axis cs:1,3) {}; \node[lab, below right] at (axis cs:1,3) {$(1,3)$};
  \node[pt] at (axis cs:-1.26,0) {}; \node[lab, below left] at (axis cs:-1.26,0) {$\left(-\sqrt[3]{2},0\right)$};
\end{axis}
\end{tikzpicture}

{\scriptsize Dashed: $y=x^2$; dotted: $y=\frac2x$; solid: the sum.}
\end{cols}

% -----------------------------------------------------------------
\subsection{Piecewise (hybrid) functions}
\label{sec:hybrid}

\begin{cols}
\begin{defn}{Hybrid functions}
A hybrid function uses different rules on different parts of its domain, e.g.
\[f(x)=\begin{cases} x^2+1, & x<1\\ 3-x, & x\ge1.\end{cases}\]
\begin{itemize}
  \item Sketch each rule \emph{only} over its own interval; show endpoints as open or closed.
  \item \textbf{Continuous} at a join $x=a$ if both pieces give the same value there (here both give $2$ at $x=1$).
  \item \textbf{Smooth} (differentiable) at the join if the gradients also match. Here the left gradient is $2x=2$ and the right gradient is $-1$, so there is a corner and $f$ is not differentiable at $x=1$ (see \S\ref{sec:cont}).
\end{itemize}
\end{defn}
\nextcol
\centering
\begin{tikzpicture}
\begin{axis}[mygraph, width=6.2cm, height=4.8cm, xmin=-2.1, xmax=3.6, ymin=-0.8, ymax=5.2, xtick={1}, ytick={2}]
  \addplot[plotA, domain=-2:1] {x^2+1};
  \addplot[plotB, domain=1:3.4] {3-x};
  \node[pt] at (axis cs:1,2) {}; \node[lab, above right] at (axis cs:1,2) {$(1,2)$ corner};
  \node[pt] at (axis cs:0,1) {}; \node[lab, right] at (axis cs:0,1) {$(0,1)$};
  \node[pt] at (axis cs:3,0) {}; \node[lab, above right] at (axis cs:3,0) {$(3,0)$};
\end{axis}
\end{tikzpicture}
\end{cols}

% -----------------------------------------------------------------
\subsection{Modelling with functions}
\label{sec:model}

\begin{keyres}{Recognising a model from a graph or a table}
\begin{center}
\small
\begin{tabularx}{\linewidth}{@{}lL@{}}
\toprule
Feature of the data or graph & Possible model \\
\midrule
Constant rate of change (equal differences in $y$ for equal steps in $x$) & Linear $y=mx+c$ \\
Single maximum or minimum, symmetric & Quadratic $y=a(x-h)^2+k$ \\
Constant \emph{ratio} for equal steps (equal percentage growth or decay) & Exponential $y=Ae^{kt}$ ($k>0$ growth, $k<0$ decay) \\
Rises or falls towards a limiting value (asymptote) & $y=c-Ae^{-kt}$ (e.g.\ heating, learning curves), $y=c+Ae^{-kt}$ (cooling) \\
Fast initial increase that keeps slowing (no asymptote) & $y=a\loge(x)+b$ or $y=ax^n$ with $0<n<1$ \\
Repeating pattern (tides, temperature, daylight) & $y=a\sin(n(t-b))+c$: $c$ = mean, $a$ = (max $-$ min)/2, period $=\frac{2\pi}{n}$ \\
Different behaviour over different intervals & Hybrid function (check continuity at the joins) \\
\bottomrule
\end{tabularx}
\end{center}
\textbf{Finding parameters:} substitute known points (simultaneous equations), use features (max, min, period, asymptote), or use CAS regression. \textbf{In context}, always state a sensible domain (e.g.\ $t\ge0$), include units and comment on whether the model is valid outside the data.
\end{keyres}

\begin{cols}
\begin{example}{circular model}
A day's temperature has a maximum of $28^\circ$C at 3 pm ($t=15$) and a minimum of $12^\circ$C at 3 am ($t=3$). Model $T=A\sin\!\left(\frac{\pi}{12}(t-b)\right)+c$.\par
Mean $c=\frac{28+12}{2}=20$, amplitude $A=\frac{28-12}{2}=8$, period 24 h so $n=\frac{2\pi}{24}=\frac{\pi}{12}$.\par
The maximum occurs when $\frac{\pi}{12}(15-b)=\frac{\pi}{2}$, so $b=9$.\par
$T(t)=8\sin\!\left(\frac{\pi}{12}(t-9)\right)+20$. Check: $T(3)=8\sin\!\left(-\frac{\pi}{2}\right)+20=12$. \tick
\end{example}
\nextcol
\begin{example}{exponential model}
A population is $500$ at $t=0$ and $2000$ at $t=3$. Model $N=Ae^{kt}$.\par
$N(0)=A=500$. \ $500e^{3k}=2000\Rightarrow e^{3k}=4\Rightarrow k=\frac{\loge4}{3}=\frac{2\loge2}{3}$.\par
So $N(t)=500e^{\frac{2\loge2}{3}t}=500\cdot4^{t/3}$.\par
The population doubles every $\frac{\loge2}{k}=1.5$ time units.
\end{example}
\end{cols}
